%! Intercepts, Zeros, and Extrema — Unit 2, Lesson 2.4 — Warm-Up
\documentclass[10pt]{article}
\usepackage{atda-article}
\usepackage{atda-boxes}
\newcommand{\TallMath}[1]{$\displaystyle #1\rule[-1.4em]{0pt}{3.2em}$}
\setlength{\workrowsep}{6pt}   % handwriting room; moves blank and key together

\begin{document}

\pageheader{Unit 2, Lesson 2.4}{Warm-Up}
% NAMESTRIP: no \namedateperiod here. The student writes their name once, on the
% cover this component is stapled behind. See references/conventions.md.

\vspace{0.15in}

\begin{notesbox}{Two Zeros, One Number Line}
\small
Five quick items. Items 1 and 4 both ask ``where is the output $0$?'' --- once from a table, once
from an equation. That question has a name, and you get it today.

\begin{enumerate}[leftmargin=1.6em, itemsep=4pt, topsep=4pt]

\item \textbf{One rule, one table.} Let $f(x)=3x-6$. Fill in every cell.
\smallskip

\renewcommand{\arraystretch}{1.3}
\begin{tabular}{@{}|c|c|c|c|c|c|@{}}
\hline
$x$ & $0$ & $1$ & $2$ & $3$ & $4$ \\ \hline
$f(x)=3x-6$ & \blank{1.0cm} & \blank{1.0cm} & \blank{1.0cm} & \blank{1.0cm} & \blank{1.0cm} \\
\hline
\end{tabular}
\smallskip

(a) Which \emph{input} gives an output of $0$? \blank{2.2cm} \quad
(b) What is the \emph{output} when the input is $0$? \blank{2.2cm}

\item \textbf{Last night's homework, in two numbers.} Item 10(b) asked for the highest profit the
yearbook staff can reach and the price at which it happens. Give \emph{both} numbers.

\writelines{1}

\item \textbf{Writing a stretch of numbers down} (Lesson 2.1). \quad (a) Write $2 < x < 7$ in
interval notation. \quad (b) Write $[0,5]$ as an inequality.

\writelines{1}

\item \textbf{From a vertex to an equation} (Lesson 2.3). A parabola has exactly the parent's shape
and its lowest point is $(4,-9)$. Write its equation, then find the two inputs where $y=0$.
\begin{work}
  y &= (x-4)^2-9 \\
  0 &= (x-4)^2-9 \\
  9 &= (x-4)^2 \\
  x-4 &= \pm 3 \\
  x &= 1 \text{ or } x = 7
\end{work}

\item \textbf{Lowest to highest.} Put these five numbers in order, lowest first.
\quad $-1$, \quad $-6$, \quad $0$, \quad $-4$, \quad $2$

\writelines{1}

\end{enumerate}
\end{notesbox}

\end{document}
